\documentclass[11pt]{article}

\usepackage[margin=0.9in]{geometry}
\usepackage[T1]{fontenc}
\usepackage{lmodern}
\usepackage{microtype}
\usepackage{xcolor}
\usepackage[hidelinks]{hyperref}
\usepackage{parskip}
\hypersetup{
	pdftitle={Response to Reviewers for Governance Requirements for Coordinating Clinical AI Agents},
	pdfauthor={Bruce Changlong Xu, Ayush Chopra, and Alexander J. Ryu}
}

\definecolor{reviewblue}{HTML}{245A78}
\newcommand{\reviewer}[1]{\section*{#1}}
\newcommand{\comment}[1]{\subsection*{\textcolor{reviewblue}{#1}}}
\newcommand{\response}{\noindent\textbf{Response.}\ }

\title{\textbf{Response to Reviewers}\\[0.5em]
\large Governance Requirements for Coordinating Clinical AI Agents}
\author{Bruce Changlong Xu, Ayush Chopra, and Alexander J. Ryu}
\date{}

\begin{document}
\maketitle

We thank the reviewers for their careful and constructive assessments. The manuscript has been
substantially restructured in response. Most importantly, we no longer present CIRC as a four-level
maturity model, a validated architecture, or a proposed standard. The revised manuscript defines CIRC
as a non-exhaustive taxonomy of clinical-agent risks and a set of independent governance levers that
can be implemented through existing standards, local infrastructure, or different architectural
patterns.

The original toy simulations, quantitative performance-style claims, oncology vignette, and
population-level clinical prediction example have been removed. CliniCARE-Bench is now used only as a
worked example of trace-observable agent-local risk. The manuscript states repeatedly that this
single-agent benchmark does not validate CIRC or test multi-agent coordination. Interaction,
workflow, institutional, and clinical-outcome claims remain an evaluation agenda rather than reported
findings.

The comments below are faithful paraphrases of the reviewer concerns rather than quotations.

\reviewer{Reviewer 1}

\comment{Comment 1: Clarify the central contribution and the relationship between individual-agent and multi-agent safety}

\response We made interaction risk the organizing problem of the manuscript. The Introduction now
distinguishes failures in an agent's own investigation from failures created by shared state,
dependencies, handoffs, finite resources, and institutional response. The revised risk taxonomy
separates agent-local, interaction-and-workflow, and system-and-institutional risks, while noting that
a single event can occupy more than one domain.

\comment{Comment 2: The four CIRC levels imply an unvalidated progression or maturity hierarchy}

\response We removed the four-level structure in full. CIRC now contains seven independent levers:
identity and attribution; scope and authorization; structured intent and shared state; dependency
coordination; resource and population awareness; operational human oversight; and monitoring and
failure containment. Deployments may select different combinations and strengths according to their
risk profile. The manuscript states that the levers are neither individually necessary nor
collectively sufficient and do not form a cumulative sequence.

\comment{Comment 3: The simulations require methodological support and should not be presented as evidence of effectiveness}

\response We removed all three toy simulations, their figures, and the associated numerical claims
rather than moving them to supplementary material. The revised paper reports no experiment purporting
to show that CIRC improves clinical or operational outcomes. It instead separates three future
evaluation questions: whether a risk is observable, whether a control detects or contains it while
preserving valid work, and whether the control improves outcomes. The last question is explicitly
reserved for prospective evaluation.

\comment{Comment 4: Ground the framework in concrete, measurable behavior}

\response We added operational definitions of interaction risk, collision, scope violation,
dependency violation, and failure containment. We also added a non-normative structured coordination
message with identity, delegated authority, patient and encounter context, state version,
dependencies, expiry, reversal, escalation, and provenance fields. Each governance lever is now linked
to illustrative observable evidence.

\reviewer{Reviewer 2}

\comment{Comment 1: CIRC is presented inconsistently as a taxonomy, maturity model, architecture, protocol, and governance framework}

\response We adopted one narrower claim throughout. CIRC is now described as an
architecture-neutral governance framework that connects a non-exhaustive risk taxonomy to candidate
controls and evidence requirements. It is not presented as a wire protocol, implemented architecture,
certification scheme, or deployment-ready standard.

\comment{Comment 2: The manuscript does not establish why the proposed controls are distinct from existing standards and infrastructure}

\response We replaced the broad insufficiency claim with two structured comparisons. The first
identifies capabilities already provided by FHIR, SMART on FHIR, CDS Hooks, A2A, MCP, and workflow and
policy engines, then separates residual semantics, authorization, workflow, governance, and
portability questions. The second maps every CIRC lever to existing substrate, the residual contract
and gap type, and illustrative evidence. We now state directly that profiles, extensions, workflow
engines, policy engines, EHR-native services, or local governance may implement many or all needed
controls without a new transport protocol.

\comment{Comment 3: The seven levers require a derivation and should not be claimed as minimal or sufficient}

\response We added the selection rationale. The seven levers are a functional decomposition of the
operational questions a deployment must answer: who acted, whether the action was permitted, what
state and intent governed it, which dependencies applied, whether shared resources were affected, who
assumed responsibility, and how failure was contained. We explicitly characterize this as a starting
point rather than an empirical derivation or closed ontology. Privacy, transport security, clinical
validation, and model performance remain cross-cutting obligations outside this decomposition.

\comment{Comment 4: Central orchestration and other architectures may address the same risks without direct agent-to-agent communication}

\response We agree. The revised manuscript compares central workflow orchestration, EHR-native
workflow control, policy and identity layers, supervisory agents, distributed agent protocols, and
human-led coordination. The architecture figure is labeled as one possible centralized realization,
not a required topology. The manuscript now treats protocol fields as most useful at organizational or
vendor boundaries and does not prefer decentralized communication over orchestration.

\comment{Comment 5: The coordination layer creates its own common-mode and security risks}

\response We added a failure model covering compromised identity, stale or replayed messages,
incorrect dependency graphs, deadlock and retry amplification, misleading resource signals,
coordination-service outage, network partition, and failure of the human control path. The revised text
requires deployment-specific trust anchors, revocation, state guarantees, and degraded behavior. It
also distinguishes safe modes for read-only, reversible administrative, consequential clinical, and
emergency actions, and specifies how control returns to a named role.

\comment{Comment 6: Human review is not an operational control unless ownership and non-response behavior are specified}

\response We replaced generic human-review language with an operational escalation contract. It
includes an accountable role and assignee, trigger and evidence, requested action, urgency and
deadline, blocked or continuing state, fallback recipient and behavior, and a closure event. The
manuscript also identifies suspension and revocation authority, separation of duties, duplicate-alert
suppression, alert budgets, and measurable assignment, acknowledgement, fallback, override, and
closure outcomes.

\comment{Comment 7: The oncology vignette and population-level prediction claim extend beyond coordination governance}

\response We removed the oncology vignette and the population-derived clinical prediction claim.
Resource and population awareness is now limited to operational concerns such as capacity, demand,
priority, concentration, fairness, and common-mode failure. Any population-level clinical prediction
is explicitly described as a separate model requiring independent development and validation.

\reviewer{Reviewer 3}

\comment{Comment 1: Define the actors, terms, and proposed information exchange more precisely}

\response The revised Introduction defines a clinical agent and distinguishes administrative,
supervised clinical, and clinically consequential agents. The operational-definitions subsection
makes the principal failure terms testable, and the structured-message example shows the information
that a receiver would need to authenticate an actor, evaluate authority, compare state, enforce
prerequisites and expiry, deduplicate retries, and route unresolved conflict. The example can map to a
FHIR profile, A2A extension, or local workflow schema and is explicitly non-normative.

\comment{Comment 2: Explain the relationship between CIRC and the Ripple Effect Protocol}

\response We now state that the Ripple Effect Protocol motivates a possible population-level failure
pattern but does not validate CIRC, establish transfer to healthcare, or define a required CIRC
mechanism. CIRC does not require REP's sensitivity-sharing approach. Aggregate demand or sensitivity
signals are described only as one candidate implementation that would require separate privacy,
manipulation, fairness, workflow, and clinical-safety evaluation.

\comment{Comment 3: Map CIRC to concrete regulatory duties and identify what the framework does not resolve}

\response We rewrote the regulatory section at the level of individual software functions,
intended use, statutory role, and deployment context. The U.S. discussion distinguishes non-device
clinical decision support from device software functions and explains the limited evidentiary role of
predetermined change control plans. The EU discussion maps Articles 9, 12, 14, 15, 17, 26, and 72 to
risk management, record keeping, oversight, robustness, quality management, deployer duties, and
post-market monitoring. A new table separates existing obligations, CIRC operational evidence, and
questions CIRC cannot decide, including classification, conformity assessment, safety and
effectiveness, intended use, liability, staffing, and local law.

\comment{Comment 4: Strengthen the scientific and standards grounding and remove unsupported claims}

\response We audited the bibliography and removed sources retained only for the deleted vignette or
superseded regulatory discussion. Every retained bibliography entry is now cited in the text. The
revised manuscript adds workflow, state-machine replication, zero-trust, sociotechnical safety,
automation, and systems-theoretic safety sources where those concepts are used. Claims of novelty,
necessity, sufficiency, validation, clinical benefit, and regulatory conformity have been narrowed or
removed.

\section*{Summary of the revised evidentiary claim}

The revised manuscript does not claim that CIRC has been empirically validated. Its contribution is a
starting vocabulary and crosswalk connecting clinical-agent risks, governance controls, existing
infrastructure, and observable evidence. CliniCARE-Bench illustrates only that some agent-local
governance properties can be assessed from traces. Multi-agent replay, system stress testing,
prospective shadow deployment, and monitored clinical use remain future work.

\end{document}
